% ============================================================
%  Greek Abstract  (Περίληψη)
% ============================================================
\chapter*{%
  \selectlanguage{greek}Εκτενής Περίληψη στα Ελληνικά\selectlanguage{english}%
}
\addcontentsline{toc}{chapter}{%
  \selectlanguage{greek}Εκτενής Περίληψη στα Ελληνικά\selectlanguage{english}%
}

\selectlanguage{greek}
\renewcommand{\thesection}{\arabic{section}}
\section{Εισαγωγή και Κίνητρο}
Τα νέφη σημείων έχουν γίνει η κυρίαρχη αναπαράσταση για τρισδιάστατα γεωμετρικά δεδομένα σε εφαρμογές όπως η αυτόνομη οδήγηση σε οχήματα με αισθητήρες LiDAR,
συστήματα καταγραφής βίντεο με όγκο για τηλεπαρουσία και εφαρμογές επαυξημένης πραγματικότητας, και ψηφιοποίηση πολιτιστικής κληρονομιάς και όλα αυτά σε ακρίβεια χιλιοστού.
Σε κάθε μία από αυτές τις εφαρμογές, ο όγκος των υπό επεξεργασία δεδομένων είναι τεράστιος, καθιστώντας την αποτελεσματική συμπίεση απαραίτητη για πρακτική χρήση.

Οι παραδοσιακοί κωδικοποιητές του MPEG, G-PCC~\cite{graziosi2020gpcc} και V-PCC~\cite{schwarz2018emerging}, αντιμετωπίζουν εν μέρει αυτήν την ανάγκη, χωρίς όμως να φτάνουν στην αποδοτικότητα που επιδεικνύουν σήμερα τα μαθησιακά συστήματα.
Μαθησιακοί κωδικοποιητές όπως ο PCGCv2~\cite{wang2021multiscale} και ο Unicorn~\cite{wang2025unicorn} έχουν πλέον εξισώσει ή υπερβεί τον G-PCC σε απόδοση ρυθμού-παραμόρφωσης, αλλά εισάγουν ένα συστηματικό γεωμετρικό σφάλμα:
την \textit{υπερλείανση} (\textit{oversmoothing}), που συγκεντρώνεται σε περιοχές υψηλής καμπυλότητας όπως ακμές, ράχες και λεπτές δομικές επιφάνειες.

Το φαινόμενο αυτό πηγάζει από μία θεμελιώδη ιδιότητα των μοντέλων εντροπίας: σπάνια, λεπτομερή μοτίβα κατοχής φέρουν υψηλό κόστος κωδικοποίησης.
Οι περιοχές υψηλής καμπυλότητας παράγουν ακριβώς τέτοια μοτίβα, οπότε ο κωδικοποιητής αποφεύγει σιωπηλά την πιστή αναπαράστασή τους.
Σε χαμηλούς ρυθμούς, όπου ο χώρος λανθάνουσας αναπαράστασης είναι περιορισμένος, η τάση αυτή εντείνεται, με αποτέλεσμα σημαντική γεωμετρική υποβάθμιση στις περιοχές ακμών και λεπτών δομών.

Το αποτέλεσμα είναι μια χωρικά ανομοιόμορφη υποβάθμιση: οι ακμές φέρουν δυσανάλογα μεγάλο μέρος του σφάλματος ενώ οι επίπεδες επιφάνειες ανακατασκευάζονται σχεδόν άρτια.
Οι τυπικές μετρικές ποιότητας, όπως η D1~PSNR~\cite{tian2017geometric}, υπολογίζουν τον μέσο όρο σε ολόκληρη την επιφάνεια αποδίδοντας ίσο βάρος σε κάθε σημείο, αποκρύπτοντας έτσι αυτή την ανισορροπία.
Ταυτόχρονα, το πρόβλημα της υπερλείανσης δεν έχει μελετηθεί σε βάθος, ούτε έχουν προταθεί μέθοδοι διόρθωσης του που να μπορούν να εφαρμοστούν σε
αποκωδικοποιημένα νέφη σημείων χωρίς τροποποίηση του υποκείμενου κωδικοποιητή.
Ένας κωδικοποιητής ο οποίος θα μπορούσε να ανακατασκευάσει μία επιφάνεια με 90\% πιστότητα ενώ συστηματικά καταστρέφει τις λεπτομέρειες στις ακμές, θα μπορούσε να
έχει υψηλό D1~PSNR αλλά να είναι ανεπαρκής για κάποιες από τις προαναφερθείσες εφαρμογές όπου η διατήρηση της γεωμετρικής λεπτομέρειας είναι κρίσιμη.

% =============================================================
\section{Συνεισφορές}
Η παρούσα διπλωματική εργασία εστιάζει σε δύο αλληλένδετα μεταξύ τους κενά. Την \emph{αδυναμία μέτρησης} της υπερλείανσης με τις υπάρχουσες μετρικές καθώς και την
\emph{έλλειψη μεθόδων διόρθωσης} που μπορούν να εφαρμοστούν σε αποκωδικοποιημένα νέφη σημείων χωρίς τροποποίηση του υποκείμενου κωδικοποιητή.
Παρουσιάζουμε, λοιπόν, δύο κύριες συνεισφορές:

\begin{enumerate}
  \item Εισάγουμε μία \emph{μετρική PSNR με διαστρωμάτωση καμπυλότητας}, η οποία κατανέμει τα σημεία βάσει της τοπικής καμπυλότητας επιφάνειας και
  αξιολογεί την ποιότητα ξεχωριστά για περιοχές ακμών και επίπεδες περιοχές, αποκαλύπτοντας υποβάθμιση που δεν μπορούν να συλλάβουν οι μετρικές που υπολογίζουν συνολικούς μέσους όρους.
  \item Προτείνουμε ένα \emph{ρυθμο-προσαρμοζόμενο αραιό U-net δίκτυο μεταεπεξεργασίας} για τη διόρθωση γεωμετρικών αλλοιώσεων σε
  αποκωδικοποιημένα νέφη σημείων χωρίς τροποποίηση του υποκείμενου κωδικοποιητή. Ένα μοντέλο 2,5 εκατομμυρίων παραμέτρων μπορεί να διαχειριστεί μόνο του πολλαπλούς ρυθμούς bit μέσω μίας κεφαλής Feature-wise Linear Modulation (FiLM)~\cite{perez2018film}, η οποία προσαρμόζει τις προβλέψεις συνεχούς μετατόπισης στον ρυθμό συμπίεσης, ακολουθώντας ένα λογαριθμικά κλιμακωμένο βαθμωτό bits-ανα-σημείο (log-BPP).
\end{enumerate}

% =============================================================
\section{Διαγνωστική ποσοτικοποίηση της υπερλείανσης και πειραματική αξιολόγηση}
Η στρωματοποιημένη κατά καμπυλότητα μετρική υπολογίζεται σε τρία στάδια: εκτίμηση καμπυλότητας, διαστρωμάτωση και αξιολόγηση PSNR ανά στρώμα.

\paragraph{Υπολογισμός καμπυλότητας.}
Για κάθε σημείο $p$ στο \emph{αρχικό} νέφος σημείων $\mathcal{P}$ (δηλαδή, το πρωτότυπο νέφος σημείων πριν από τη συμπίεση),
υπολογίζουμε την τοπική καμπυλότητα $\kappa(p)$ από τις ιδιοτιμές του τοπικού πίνακα συνδιακύμανσης που κατασκευάζεται από τους $k=30$ πλησιέστερους γείτονες.
Η μικρότερη ιδιοτιμή, κανονικοποιημένη από το άθροισμα και των τριών, δίνει τη διακύμανση της τοπικής επιφάνειας:
τιμή κοντά στο 0 αντιστοιχεί σε επίπεδη περιοχή, ενώ τιμή κοντά στο $1/3$ σε ισότροπα κατανεμημένη γειτονιά.

\paragraph{Διαστρωμάτωση.}
Τα σημεία διαχωρίζονται σε τρία στρώματα με βάση τα τριτημόρια της κατανομής καμπυλότητας, δημιουργώντας έτσι τα στρώματα ακμών, μέσης καμπυλότητας και
επίπεδων περιοχών. Χρησιμοποιούμε τα 67ο και 33ο ποσοστημόριο ως κατώφλια για να διαχωρίσουμε τα στρώματα, εξασφαλίζοντας ότι κάθε στρώμα περιέχει το
ένα τρίτο των σημείων του αρχικού νέφους σημείων. Αυτή η προσαρμοστική διαστρωμάτωση εξασφαλίζει ότι όλα τα στρώματα είναι ισοδύναμα σε μέγεθος,
καθιστώντας την προσέγγιση ανεξάρτητη από το σύνολο δεδομένων και επιτρέποντας δίκαιη σύγκριση μεταξύ διαφορετικών συνόλων δεδομένων και κωδικοποιητών.

\paragraph{Αξιολόγηση PSNR ανά στρώμα.}
Για κάθε στρώμα, υπολογίζουμε το PSNR χρησιμοποιώντας την αντίστροφη κατεύθυνση (από το πρωτότυπο στο ανακατασκευασμένο νέφος σημείων) για να μετρήσουμε
πόσο καλά η ανακατασκευή καλύπτει την αρχική επιφάνεια, αντί να μετράμε εαν κάθε σημείο της ανακατασκευασμένης επιφάνειας είναι κοντά σε κάποιο πρωτότυπο.
Η κατευθυντικότητα της μετρικής δεν είναι απλά μία σχεδιαστική επιλογή αλλά είναι κρίσιμη για την αποκάλυψη της υποβάθμισης στις περιοχές ακμών.
Σε χαμηλούς ρυθμούς συμπίεσης, η καμπυλότητα του \emph{ανακατασκευασμένου} νέφους, παρουσιάζει αρνητική συσχέτιση με την καμπυλότητα του πρωτότυπου νέφους σε περιοχές
ακμών (συντελεστής συσχέτισης Pearson $r=-0.33$), υποδεικνύοντας ότι οι περιοχές υψηλής καμπυλότητας τείνουν να υποβαθμίζονται περισσότερο λόγω της υπερλείανσης.
Χρησιμοποιώντας την καμπυλότητα του αποκωδικοποιημένου νέφους για διαστρωμάτωση θα κατέτασσε τα σημεία που έχουν υποστεί την περισσότερη υπερλείανση ως χαμηλής καμπυλότητας,
αποκρύπτοντας έτσι την υποβάθμιση που προσπαθούμε να μετρήσουμε.

% =============================================================
\section{Διερεύνηση της Υπερλείανσης στον κωδικοποιητή PCGCv2}

Εφαρμόζουμε τη μετρική στον PCGCv2 σε τέσσερις ακολουθίες του συνόλου δεδομένων
8i~Voxelized Full Bodies (8iVFB)~\cite{8ivfb}: \textit{longdress}, \textit{loot},
\textit{redandblack} και \textit{soldier}.
Οι ακολουθίες αυτές αποτελούν το τυπικό σύνολο αξιολόγησης της MPEG για
συμπίεση νεφών σημείων και καλύπτουν ένα ευρύ φάσμα γεωμετρικής πολυπλοκότητας.

Τα αποτελέσματα αποκαλύπτουν τρεις συστηματικές ιδιότητες.

\paragraph{Συστηματική χωρική μεροληψία.}
Η διαφορά PSNR μεταξύ επίπεδων περιοχών και ακμών
$\Delta_{\mathrm{smooth}} = \mathrm{PSNR}_{\mathrm{flat}} - \mathrm{PSNR}_{\mathrm{edge}}$
είναι αυστηρά θετική σε όλους τους 28 συνδυασμούς ακολουθίας και ρυθμού που αξιολογήθηκαν
(επτά σημεία ρυθμού, τέσσερις ακολουθίες).
Ο PCGCv2 ανακατασκευάζει \emph{πάντα} τις επίπεδες περιοχές καλύτερα από τις ακμές,
χωρίς καμία εξαίρεση.

\paragraph{Το χάσμα διευρύνεται με τον ρυθμό.}
Παραδόξως, το χάσμα υπερλείανσης \emph{αυξάνεται} καθώς αυξάνεται ο ρυθμός μετάδοσης.
Στον υψηλότερο αξιολογούμενο ρυθμό (r7), το χάσμα είναι 2 έως 3 φορές μεγαλύτερο
από ό,τι στον χαμηλότερο (r1) για κάθε ακολουθία, κυμαινόμενο από 2,5 έως 4,5~dB.
Αυτό αποδεικνύει ότι η υπερλείανση δεν είναι απλώς ένα αποτέλεσμα λόγω χαμηλού ρυθμού συμπίεσης
που εξαφανίζεται με περισσότερα bits: ο κωδικοποιητής κατανέμει τα επιπλέον bits
κατά προτίμηση στις επίπεδες επιφάνειες, διευρύνοντας το χάσμα αντί να το κλείσει.

\paragraph{Εξάρτηση από ακολουθία.}
Το μέγεθος του χάσματος ποικίλλει σημαντικά μεταξύ ακολουθιών: η \textit{redandblack}
εμφανίζει το μεγαλύτερο χάσμα (2,4 έως 4,5~dB), συνεπές με τη λεπτή υφή της επιφάνειάς της
(πρόσωπο, μαλλιά, ύφασμα), ενώ η \textit{loot} εμφανίζει το μικρότερο χάσμα σε χαμηλούς
ρυθμούς (0,2~dB στο r1).

Η εφαρμογή της ίδιας αξιολόγησης στον G-PCC αποκαλύπτει ένα ποιοτικά διαφορετικό
προφίλ αλλοιώσεως: το χάσμα κορυφώνεται σε ενδιάμεσους ρυθμούς (μέση τιμή 6,4~dB)
και στη συνέχεια \emph{μειώνεται} σε υψηλότερους ρυθμούς, σχηματίζοντας μη μονοτονική
καμπύλη.
Αυτή η αντίθεση επιβεβαιώνει ότι το μονοτόνως αυξανόμενο χάσμα είναι δομική ιδιότητα
του παραδείγματος κωδικοποίησης κατοχής που ακολουθούν οι μαθησιακοί κωδικοποιητές,
και όχι της γεωμετρικής συμπίεσης γενικά.

\paragraph{Σύγκλιση σε τετριμμένη λύση: επιπτώσεις στο σχεδιασμό μεταεπεξεργασίας.}
Η εξέταση της κατανομής των μεγεθών της επιδιωκόμενης μετατόπισης (διανύσματα
πλησιέστερου γείτονα από αποκωδικοποιημένα σε αρχικά σημεία) αποκαλύπτει ότι στον
ρυθμό r2, το 60\% των αποκωδικοποιημένων σημείων ταυτίζεται ήδη με αρχικά voxels
και δεν χρειάζεται διόρθωση.
Στον ρυθμό r7, το ποσοστό αυτό ανέρχεται στο 89\%.
Αυτή η ιδιότητα \emph{κυριαρχίας μηδενικών} έχει κρίσιμη συνέπεια: ένα αφελές
μοντέλο παλινδρόμησης ελαχιστοποιεί την αναμενόμενη L1 απώλεια προβλέποντας μηδέν
για κάθε σημείο, καταρρέοντας σε μία τετριμμένη λύση που αγνοεί τα σημεία εκεί όπου
η διόρθωση είναι πιο απαραίτητη.

% =============================================================
\section{Δίκτυο Μεταεπεξεργασίας}

\subsection{Διατύπωση Προβλήματος}

Έστω $\hat{\mathcal{P}} = \{\hat{p}_i\}_{i=1}^{N} \subset \mathbb{Z}^3$ το
αποκωδικοποιημένο νέφος σε ακέραιες συντεταγμένες voxel (ανάλυση $G=1024$) και
$\mathcal{P}^* = \{p^*_j\}_{j=1}^{M}$ το αρχικό νέφος αναφοράς.
Το δίκτυο μεταεπεξεργασίας προβλέπει μια τρισδιάστατη μετατόπιση
$\Delta_i \in \mathbb{R}^3$ για κάθε αποκωδικοποιημένο σημείο:
\[
  \hat{\mathcal{P}}_{\mathrm{refined}}
    = \bigl\{\hat{p}_i + \Delta_i\bigr\}_{i=1}^{N},
  \quad
  \Delta_i = f_\theta(\hat{\mathcal{P}},\,r)_i,
  \quad r = \log_2(\mathrm{bpp}).
\]
Η έξοδος έχει τον ίδιο αριθμό σημείων με το αποκωδικοποιημένο νέφος.
Οι μετατοπίσεις περιορίζονται σε ${\pm}5$ voxels μέσω ενεργοποίησης $\tanh$,
καλύπτοντας όλες τις ρεαλιστικές διορθώσεις με περιθώριο.
Τα χαρακτηριστικά εισόδου είναι οι τρεις κανονικοποιημένες συντεταγμένες voxel
σε $[-1,1]^3$· πείραμα αφαίρεσης επιβεβαιώνει ότι η προσθήκη εκτιμώμενης
καμπυλότητας ως χαρακτηριστικό εισόδου \emph{επιδεινώνει} την απόδοση
(συζητείται παρακάτω).

\subsection{Αρχιτεκτονική: Αραιό U-Net Τριών Επιπέδων}

Ο κεντρικός κορμός είναι ένα αραιό U-Net τριών επιπέδων που λειτουργεί άμεσα
στο πλέγμα κατοχής του αποκωδικοποιημένου νέφους, υλοποιημένο με τη βιβλιοθήκη
MinkowskiEngine~\cite{choy2019minkowski}.
Κάθε επίπεδο κωδικοποιητή υποδιπλασιάζει την ανάλυση και διπλασιάζει τον αριθμό
καναλιών (32 $\to$ 64 $\to$ 128)· ο αποκωδικοποιητής αντιστρέφει αυτή τη διαδικασία
με συνδέσεις παράκαμψης (skip connections).
Η πολυ-κλιμακωτή δομή παρέχει ταυτόχρονα ευαισθησία στη τοπική γεωμετρική
λεπτομέρεια (ακμές) και ευρύτερο γεωμετρικό πλαίσιο (τοπολογία επιφάνειας),
κάτι που δεν μπορεί να επιτύχει ένα MLP ανεξάρτητο σημείων ή μια ολική λειτουργία.

\subsection{Κεφαλή FiLM για Προσαρμογή Ρυθμού}

Ένα διπλόστρωτο MLP απεικονίζει το βαθμωτό $r = \log_2(\mathrm{bpp})$ σε παραμέτρους
FiLM 64 διαστάσεων, χωρισμένες σε 32 τιμές κλίμακας $\boldsymbol{\gamma}$ και 32 τιμές
μετατόπισης $\boldsymbol{\beta}$, οι οποίες τροποποιούν το διάνυσμα χαρακτηριστικών
κάθε ενεργού voxel:
\[
  \mathbf{h}'_i = \boldsymbol{\gamma}(r) \odot \mathbf{h}_i + \boldsymbol{\beta}(r).
\]
Η κωδικοποίηση log-BPP γραμμικοποιεί τη σχέση ρυθμού-παραμόρφωσης, ώστε ίσα
διαστήματα σε αυτή την κλίμακα να αντιστοιχούν σε κατά προσέγγιση ίσες αλλαγές
παραμόρφωσης.
Το MLP της κεφαλής FiLM προσθέτει μόλις 4.224 παραμέτρους (0,17\% του συνολικού
μεγέθους μοντέλου), επιτρέποντας σε ένα μοντέλο 2,5 εκατομμυρίων παραμέτρων να
εξυπηρετεί ταυτόχρονα όλα τα σημεία λειτουργίας ρυθμού-παραμόρφωσης.

\subsection{Συνάρτηση Απώλειας: Δυναμική Απόσταση Chamfer}

Οι τυπικές συναρτήσεις απώλειας παλινδρόμησης μετατόπισης (L1, smooth-L1,
L1 σταθμισμένη με καμπυλότητα) αποτυγχάνουν λόγω δύο δομικών προβλημάτων.

\emph{Κατάρρευση λόγω κυριαρχίας μηδενικών}: στον ρυθμό r2, η διαμέσιος
των επιδιωκόμενων τιμών ανά συντεταγμένη είναι μηδέν, οπότε ο βέλτιστος σταθερός
προβλέπτης υπό L1 είναι μηδέν.
Κάθε δίκτυο εκπαιδευμένο με αυτή την απώλεια καταρρέει σε προβλέψεις κοντά στο
μηδέν εντός 5 έως 10 εποχών, γεγονός που επιβεβαιώθηκε εμπειρικά για όλες τις
τυπικές παραλλαγές απώλειας.

\emph{Θόρυβος πλησιέστερου γείτονα}: σε περιοχές ακμών, μια μικρή μεταβολή της
θέσης ενός αποκωδικοποιημένου σημείου μπορεί να αλλάξει τον πλησιέστερο γείτονά
του από τη μία πλευρά της ακμής στην άλλη, προκαλώντας ασυνεχή άλμα στον στόχο.
Αυτή η αστάθεια είναι εντονότερη ακριβώς εκεί που η διόρθωση είναι πιο απαραίτητη.

Επιλύουμε και τα δύο προβλήματα εγκαταλείποντας πλήρως τους προϋπολογισμένους
ανά σημείο στόχους.
Αντ' αυτού, ελαχιστοποιούμε άμεσα τη συμμετρική Chamfer Distance μεταξύ του
επεξεργασμένου τμήματος και της αντίστοιχης τομής του αρχικού νέφους:
\[
  \mathcal{L}_{\mathrm{CD}} =
    \frac{1}{|\hat{\mathcal{P}}_{\mathrm{ref}}|}
    \sum_{\hat{p} \in \hat{\mathcal{P}}_{\mathrm{ref}}}
      \min_{p^* \in \mathcal{P}^*_{\mathrm{crop}}} \|\hat{p} - p^*\|^2
    +
    \frac{1}{|\mathcal{P}^*_{\mathrm{crop}}|}
    \sum_{p^* \in \mathcal{P}^*_{\mathrm{crop}}}
      \min_{\hat{p} \in \hat{\mathcal{P}}_{\mathrm{ref}}} \|p^* - \hat{p}\|^2.
\]
Το πρώτο σκέλος (εμπρός) τιμωρεί τα επεξεργασμένα σημεία που βρίσκονται μακριά
από την αρχική επιφάνεια· το δεύτερο (ανάστροφο) τιμωρεί τις περιοχές της αρχικής
επιφάνειας που δεν καλύπτονται από το επεξεργασμένο νέφος.
Η αντιστοίχιση πλησιέστερου γείτονα επαναϋπολογίζεται σε κάθε βήμα εκπαίδευσης
βάσει των τρεχουσών επεξεργασμένων θέσεων, εξ ου και ο χαρακτηρισμός «δυναμική».
Για σημεία που ήδη βρίσκονται σωστά τοποθετημένα, η συνεισφορά στην κλίση είναι
ελάχιστη ανεξαρτήτως ποσοστού μηδενικών, εξαλείφοντας έτσι τη κυριαρχία μηδενικών.
Ο θόρυβος πλησιέστερου γείτονα επίσης εξαλείφεται, διότι το δίκτυο μπορεί να
μετακινεί σημεία ελεύθερα εντός $[-5,5]^3$ voxels και η αντιστοίχιση προσαρμόζεται
αναλόγως.

Μία κρίσιμη λεπτομέρεια υλοποίησης είναι ότι η τομή του αρχικού νέφους πρέπει να
χρησιμοποιεί \emph{μηδενικό} περιθώριο (zero padding).
Με μη μηδενικό περιθώριο, το ανάστροφο σκέλος κυριαρχείται από αρχικά σημεία εκτός
ορίων και καταρρέει, δημιουργώντας εκφυλισμένη κλίση που ωθεί τα αποκωδικοποιημένα
σημεία προς τα όρια του τμήματος.
Με μηδενικό περιθώριο, τα δύο σκέλη ισορροπούν και το μοντέλο μαθαίνει σωστές
γεωμετρικές διορθώσεις.

% =============================================================
\section{Πειραματικά Αποτελέσματα}

\subsection{Κύρια Αποτελέσματα στο 8iVFB}

Αξιολογούμε το τελικό μοντέλο σε 400 πλαίσια και τις τέσσερις ακολουθίες του 8iVFB
σε τρία σημεία ρυθμού (r2, r4, r6).
Οι μέσες βελτιώσεις PSNR έναντι του επιπέδου αναφοράς PCGCv2 είναι:

\begin{center}
\begin{tabular}{lcc}
  \toprule
  Ρυθμός & PSNR ακμών $\delta$ (dB) & PSNR επίπεδων $\delta$ (dB) \\
  \midrule
  r2 ($\approx$0,05 bpp) & $\mathbf{+2{,}06}$ & $+1{,}47$ \\
  r4 ($\approx$0,15 bpp) & $+0{,}58$ & $+0{,}62$ \\
  r6 ($\approx$0,32 bpp) & $+0{,}27$ & $+0{,}18$ \\
  \bottomrule
\end{tabular}
\end{center}

Στον ρυθμό r2, το επεξεργασμένο PSNR ακμών (62,11~dB) πλησιάζει σημαντικά το PSNR
επίπεδων περιοχών του επιπέδου αναφοράς (62,70~dB), σχεδόν κλείνοντας το χάσμα
στρωμάτωσης.
Το Bj{\o}ntegaard Delta PSNR ολοκληρωμένο στο εύρος r2 έως r6 είναι \textbf{+0,97~dB}
στη μετρική ακμών.
Οι βελτιώσεις είναι συνεπείς σε όλες τις τέσσερις ακολουθίες, συμπεριλαμβανομένων
τριών που δεν εμφανίστηκαν στην εκπαίδευση (εύρος βελτίωσης PSNR ακμών στο r2:
1,78 έως 2,21~dB).

\subsection{Πείραμα Αφαίρεσης: Καμπυλότητα ως Χαρακτηριστικό Εισόδου}

Η αφαίρεση της εκτιμώμενης καμπυλότητας από τα χαρακτηριστικά εισόδου
\emph{βελτιώνει} τα αποτελέσματα σε όλους τους ρυθμούς.
Στο r2, η βελτίωση PSNR ακμών αυξάνεται από +1,62~dB (με καμπυλότητα) σε
+2,06~dB (χωρίς), κέρδος 0,44~dB.
Δύο παράγοντες εξηγούν αυτό: (α) η καμπυλότητα εκτιμώμενη στο αποκωδικοποιημένο
νέφος είναι πιο αναξιόπιστη ακριβώς στις περιοχές υψηλής καμπυλότητας όπου
χρειάζονται οι μεγαλύτερες διορθώσεις· (β) το πεδίο υποδοχής του δικτύου
($13^3$ voxels σε πλήρη ανάλυση) συλλαμβάνει ήδη σιωπηλά τη γεωμετρία της
τοπικής επιφάνειας από τα χαρακτηριστικά κανονικοποιημένης θέσης.
Αυτό το εύρημα έχει μεθοδολογική συνέπεια: η παροχή φαινομενικά σχετικών
χαρακτηριστικών μπορεί να επιβαρύνει τη μάθηση όταν τα χαρακτηριστικά αυτά
είναι πιο θορυβώδη ακριβώς εκεί που το σήμα έχει τη μεγαλύτερη σημασία.

% =============================================================
\section{Γενίκευση: Μεταφορά Zero-Shot}

\subsection{Άγνωστος Κωδικοποιητής: Unicorn}

Εφαρμόζουμε το μοντέλο αμεταβλήτως στον Unicorn~\cite{wang2025unicorn},
έναν κωδικοποιητή αιχμής που δεν εμφανίστηκε κατά την εκπαίδευση.
Τα αποτελέσματα αποκαλύπτουν ένα σαφές κατώφλι ρυθμού.
Σε μέτριους ρυθμούς (Unicorn R2 έως R4, bpp $> 0{,}05$), τα μοτίβα παραμόρφωσης
διαφέρουν από αυτά του PCGCv2, προκαλώντας ελαφρά υποβάθμιση ακμών
($-$0,01 έως $-$0,72~dB).
Σε ακραίους ρυθμούς (R5 έως R6, bpp $< 0{,}025$), η υπερλείανση είναι τόσο έντονη
ώστε να είναι κοινή ανεξαρτήτως κωδικοποιητή και οι διορθώσεις μεταφέρονται:
\textbf{+1,09~dB} στο R5 και \textbf{+1,36~dB} στο R6, κατά μέσο όρο σε όλες τις
ακολουθίες.

\subsection{Άγνωστο Σύνολο Δεδομένων: MVUB}

Αξιολογούμε το μοντέλο στο σύνολο δεδομένων MVUB~\cite{loop2016mvub} (καταγραφές
άνω σώματος σε ανάλυση $G=512$, ακολουθίες \textit{david} και \textit{ricardo}),
που αποτελεί δοκιμή εκτός κατανομής ως προς ανάλυση και σύνολο δεδομένων.
Χωρίς επανεκπαίδευση, το μοντέλο επιτυγχάνει θετική βελτίωση PSNR ακμών
σε χαμηλούς ρυθμούς: \textit{david} +0,92~dB και \textit{ricardo} +1,03~dB στο r2.
Σε μέτριους έως υψηλούς ρυθμούς (r4, r6), η μεταφορά είναι ουδέτερη έως αρνητική,
αντικατοπτρίζοντας το πρότυπο του Unicorn.

Συνολικά, αυτά τα πειράματα αποδεικνύουν ότι η γεωμετρική υπερλείανση χαμηλού
ρυθμού αποτελεί χαρακτηριστική αποτυχία πολλαπλών κωδικοποιητών και συνόλων
δεδομένων, και ότι ένα μοντέλο εκπαιδευμένο σε έναν κωδικοποιητή μπορεί να
γενικεύσει τις διορθώσεις του όταν η υπερλείανση είναι αρκετά σοβαρή ώστε να
κυριαρχεί στο μοτίβο παραμόρφωσης.

% =============================================================
\section{Συμπεράσματα}

Η παρούσα διπλωματική εργασία αντιμετωπίζει ένα συγκεκριμένο και υποτιμημένο
πρόβλημα στη συμπίεση νεφών σημείων με μηχανική μάθηση: τη γεωμετρική υπερλείανση
που συγκεντρώνεται στις περιοχές υψηλής καμπυλότητας και είναι αόρατη στις
συνολικές μετρικές.

Εισάγουμε μία \emph{μετρική PSNR με διαστρωμάτωση καμπυλότητας} που καθιστά
αυτό το κρυμμένο πρόβλημα ορατό, και αποδεικνύουμε ότι στον PCGCv2 το χάσμα
υπερλείανσης δεν είναι παραμόρφωση λόγω χαμηλού ρυθμού αλλά δομική ιδιότητα του
μαθησιακού παραδείγματος κωδικοποίησης κατοχής που αυξάνεται μονότονα με τον ρυθμό.
Η αντίθεση με τον G-PCC, που εμφανίζει μη μονοτονικό, μειούμενο με τον ρυθμό
χάσμα, επιβεβαιώνει ότι η υπερλείανση είναι εγγενής στο μοντέλο εντροπίας
πλέγματος κατοχής.

Προτείνουμε ένα \emph{ρυθμο-προσαρμοζόμενο αραιό U-Net δίκτυο μεταεπεξεργασίας}
με κεφαλή FiLM, εκπαιδευμένο με δυναμική Chamfer Distance απώλεια που παρακάμπτει
τους δύο δομικούς τρόπους αποτυχίας της τυπικής παλινδρόμησης μετατόπισης.
Το μοντέλο επιτυγχάνει σημαντική ανάκτηση PSNR ακμών σε χαμηλούς ρυθμούς
(+2,06~dB στο r2) ενώ βελτιώνει και τις επίπεδες περιοχές, και γενικεύεται χωρίς
επανεκπαίδευση σε άγνωστο κωδικοποιητή και άγνωστο σύνολο δεδομένων σε ακραίους
ρυθμούς συμπίεσης.

\vspace{1em}
\noindent\textbf{Λέξεις-κλειδιά:}
συμπίεση νέφους σημείων, μεταεπεξεργασία, υπερλείανση,
καμπυλότητα, αραιές συνελίξεις, FiLM, ρυθμο-προσαρμογή, PCGCv2.

% Reset citation counter so Chapter 1 starts from [1]
\setcounter{NAT@ctr}{0}

\selectlanguage{english}
\cleardoublepage
